% GENERATED FILE. Edit canonical lessons, then run npm run generate:books.
% canonical-source-hash: fnv1a64:629fe31f76d56e3a
% canonical-lessons: TA-C296-kavalturai, TA-C296-kaitu-cey, TA-C296-currivalai, TA-C296-kaiti, TA-C296-oppukkol

\chapter{The Police Close In}
\label{ch:ta-the-police-close-in}
\hlchaptermodality{296}

\begin{chapteropening}
\textbf{By the end of this chapter:} \emph{I can say what the police did.}

Complete the last lesson of chapter 296: I can say what the police did.
\end{chapteropening}

\section[\texorpdfstring{\ta{காவல்துறை} (kāvaltuṟai)}{காவல்துறை (kāvaltuṟai)}]{\ta{காவல்துறை} (kāvaltuṟai) — the police}

\label{lesson:TA-C296-kavalturai}



\begin{quote}
Before the new one: say the Tamil for a crime, then the Tamil for a culprit.
\end{quote}

\subsection*{You'll want to know: \ta{காவல்துறை}}
\textbf{\ta{காவல்துறை}} — \emph{kāvaltuṟai} — \textquotedblleft{}the police\textquotedblright{}.

\begin{quote}
\ta{திருடனைப்} \ta{பிடிப்பது} \ta{காவல்துறையைப்} \ta{பொறுத்தது}.
\end{quote}

\begin{quote}
\emph{tiruṭaṉaip piṭippatu kāvaltuṟaiyaip poṟuttatu.}
\end{quote}

\begin{quote}
Catching the thief depends on the police.
\end{quote}

Its officers are \ta{காவல்துறையினர்}; they work from a \ta{காவல்} \ta{நிலையம்}, a police station.

\begin{cousinweb}
\textbf{\ta{காவல்துறை}} (\emph{kāvaltuṟai}) is a Tamil compound built from \ta{காவல்}, \textquotedblleft{}guarding, protection\textquotedblright{} (DEDR 1416), and \ta{துறை}, \textquotedblleft{}place, branch, department\textquotedblright{} (DEDR 3370).
\end{cousinweb}

\begin{grammarlens}[title={What you've built}]
One more word for telling what happened.
\end{grammarlens}

\subsection*{Your turn}
\begin{itemize}
\raggedright
  \item \emph{Say it:} \emph{kāvaltuṟai}
  \item \emph{Say it:} \emph{kāvaltuṟai}, once more
  \item \emph{Say it:} \emph{tiruṭaṉaip piṭippatu kāvaltuṟaiyaip poṟuttatu.} — Catching the thief depends on the police.
  \item \emph{Recall:} say the Tamil for a crime, then the Tamil for a culprit, then say \emph{kāvaltuṟai} again
\end{itemize}

\begin{tcolorbox}[breakable,colback=teal!4,colframe=teal!35!black,title={Before you move on}]
What does \ta{காவல்துறை} mean? (\textquotedblleft{}The police\textquotedblright{}.) Say it once more.
\end{tcolorbox}
\section[\texorpdfstring{\ta{கைது} \ta{செய்} (kaitu cey)}{கைது செய் (kaitu cey)}]{\ta{கைது} \ta{செய்} (kaitu cey) — to arrest}

\label{lesson:TA-C296-kaitu-cey}



\begin{quote}
Before the new one: say the Tamil for a culprit, then the Tamil for the police.
\end{quote}

\subsection*{You'll want to know: \ta{கைது} \ta{செய்}}
\textbf{\ta{கைது} \ta{செய்}} — \emph{kaitu cey} — \textquotedblleft{}to arrest\textquotedblright{}.

\begin{quote}
\ta{காவலர்கள்} \ta{இரண்டு} \ta{பேரைக்} \ta{கைது} \ta{செய்தார்கள்}.
\end{quote}

\begin{quote}
\emph{kāvalarkaḷ iraṇṭu pēraik kaitu ceytārkaḷ.}
\end{quote}

\begin{quote}
The police arrested two people.
\end{quote}

Past (I …): \ta{கைது} \ta{செய்தேன்}, \emph{kaitu ceytēṉ}.

The person arrested takes -\ta{ஐ}; on its own, \ta{கைது} is an arrest.

\begin{cousinweb}
\textbf{\ta{கைது} \ta{செய்}} (\emph{kaitu cey}) joins \ta{கைது}, \textquotedblleft{}arrest\textquotedblright{}, and \ta{செய்}, \textquotedblleft{}to do\textquotedblright{}. \ta{கைது} is usually traced to Arabic \emph{qayd}, \textquotedblleft{}a bond, captivity\textquotedblright{}, which reached Tamil through Hindustani.
\end{cousinweb}

\begin{grammarlens}[title={What you've built}]
One more word for telling what happened.
\end{grammarlens}

\subsection*{Your turn}
\begin{itemize}
\raggedright
  \item \emph{Say it:} \emph{kaitu cey}
  \item \emph{Say it:} \emph{kaitu cey}, once more
  \item \emph{Say it:} \emph{kāvalarkaḷ iraṇṭu pēraik kaitu ceytārkaḷ.} — The police arrested two people.
  \item \emph{Recall:} say the Tamil for a culprit, then the Tamil for the police, then say \emph{kaitu cey} again
\end{itemize}

\begin{tcolorbox}[breakable,colback=teal!4,colframe=teal!35!black,title={Before you move on}]
What does \ta{கைது} \ta{செய்} mean? (\textquotedblleft{}To arrest\textquotedblright{}.) Say it once more.
\end{tcolorbox}
\section[\texorpdfstring{\ta{சுற்றிவளை} (cuṟṟivaḷai)}{சுற்றிவளை (cuṟṟivaḷai)}]{\ta{சுற்றிவளை} (cuṟṟivaḷai) — to cordon off; to surround}

\label{lesson:TA-C296-currivalai}



\begin{quote}
Before the new one: say the Tamil for the police, then the Tamil for to arrest.
\end{quote}

\subsection*{You'll want to know: \ta{சுற்றிவளை}}
\textbf{\ta{சுற்றிவளை}} — \emph{cuṟṟivaḷai} — \textquotedblleft{}to cordon off; to surround\textquotedblright{}.

\begin{quote}
\ta{காவலர்கள்} \ta{அந்த} \ta{வீட்டைச்} \ta{சுற்றிவளைத்தார்கள்}.
\end{quote}

\begin{quote}
\emph{kāvalarkaḷ anta vīṭṭaic cuṟṟivaḷaittārkaḷ.}
\end{quote}

\begin{quote}
The police surrounded that house.
\end{quote}

Past (I …): \ta{சுற்றிவளைத்தேன்}, \emph{cuṟṟivaḷaittēṉ}.

What is surrounded takes -\ta{ஐ}: \ta{சுற்றி}, around, and \ta{வளை}, to encircle.

\begin{cousinweb}
\textbf{\ta{சுற்றிவளை}} (\emph{cuṟṟivaḷai}) is a Tamil compound built from \ta{சுற்று}, \textquotedblleft{}to go round\textquotedblright{} (DEDR 2715), and \ta{வளை}, \textquotedblleft{}to surround, encircle\textquotedblright{} (DEDR 5313).
\end{cousinweb}

\begin{grammarlens}[title={What you've built}]
One more word for telling what happened.
\end{grammarlens}

\subsection*{Your turn}
\begin{itemize}
\raggedright
  \item \emph{Say it:} \emph{cuṟṟivaḷai}
  \item \emph{Say it:} \emph{cuṟṟivaḷai}, once more
  \item \emph{Say it:} \emph{kāvalarkaḷ anta vīṭṭaic cuṟṟivaḷaittārkaḷ.} — The police surrounded that house.
  \item \emph{Recall:} say the Tamil for the police, then the Tamil for to arrest, then say \emph{cuṟṟivaḷai} again
\end{itemize}

\begin{tcolorbox}[breakable,colback=teal!4,colframe=teal!35!black,title={Before you move on}]
What does \ta{சுற்றிவளை} mean? (\textquotedblleft{}To cordon off; to surround\textquotedblright{}.) Say it once more.
\end{tcolorbox}
\section[\texorpdfstring{\ta{கைதி} (kaiti)}{கைதி (kaiti)}]{\ta{கைதி} (kaiti) — a prisoner}

\label{lesson:TA-C296-kaiti}



\begin{quote}
Before the new one: say the Tamil for to arrest, then the Tamil for to cordon off.
\end{quote}

\subsection*{You'll want to know: \ta{கைதி}}
\textbf{\ta{கைதி}} — \emph{kaiti} — \textquotedblleft{}a prisoner\textquotedblright{}.

\begin{quote}
\ta{கைதி} \ta{இரவில்} \ta{சிறையிலிருந்து} \ta{தப்பினான்}.
\end{quote}

\begin{quote}
\emph{kaiti iravil ciṟaiyiliruntu tappiṉāṉ.}
\end{quote}

\begin{quote}
The prisoner escaped from jail at night.
\end{quote}

Plural \ta{கைதிகள்}; it goes with \ta{கைது}, an arrest.

\begin{cousinweb}
\textbf{\ta{கைதி}} (\emph{kaiti}) goes with \ta{கைது}, \textquotedblleft{}arrest\textquotedblright{}: it is usually traced to Arabic \emph{qaydī}, \textquotedblleft{}a captive\textquotedblright{}, through Hindustani.
\end{cousinweb}

\begin{grammarlens}[title={What you've built}]
One more word for telling what happened.
\end{grammarlens}

\subsection*{Your turn}
\begin{itemize}
\raggedright
  \item \emph{Say it:} \emph{kaiti}
  \item \emph{Say it:} \emph{kaiti}, once more
  \item \emph{Say it:} \emph{kaiti iravil ciṟaiyiliruntu tappiṉāṉ.} — The prisoner escaped from jail at night.
  \item \emph{Recall:} say the Tamil for to arrest, then the Tamil for to cordon off, then say \emph{kaiti} again
\end{itemize}

\begin{tcolorbox}[breakable,colback=teal!4,colframe=teal!35!black,title={Before you move on}]
What does \ta{கைதி} mean? (\textquotedblleft{}A prisoner\textquotedblright{}.) Say it once more.
\end{tcolorbox}
\section[\texorpdfstring{\ta{ஒப்புக்கொள்} (oppukkoḷ)}{ஒப்புக்கொள் (oppukkoḷ)}]{\ta{ஒப்புக்கொள்} (oppukkoḷ) — to confess; to admit}

\label{lesson:TA-C296-oppukkol}



\begin{quote}
Before the new one: say the Tamil for to cordon off, then the Tamil for a prisoner.
\end{quote}

\subsection*{You'll want to know: \ta{ஒப்புக்கொள்}}
\textbf{\ta{ஒப்புக்கொள்}} — \emph{oppukkoḷ} — \textquotedblleft{}to confess; to admit\textquotedblright{}.

\begin{quote}
\ta{கடைசியில்} \ta{அவன்} \ta{தன்} \ta{குற்றத்தை} \ta{ஒப்புக்கொண்டான்}.
\end{quote}

\begin{quote}
\emph{kaṭaiciyil avaṉ taṉ kuṟṟattai oppukkoṇṭāṉ.}
\end{quote}

\begin{quote}
In the end he confessed his crime.
\end{quote}

Past (I …): \ta{ஒப்புக்கொண்டேன்}, \emph{oppukkoṇṭēṉ}.

What is admitted takes -\ta{ஐ}. Also to agree to do something: \ta{வேலையை} \ta{ஒப்புக்கொண்டேன்}.

\begin{cousinweb}
\textbf{\ta{ஒப்புக்கொள்}} (\emph{oppukkoḷ}) is a Tamil compound built from \ta{ஒப்பு}, \textquotedblleft{}agreement, assent\textquotedblright{} (DEDR 924), and \ta{கொள்}, \textquotedblleft{}to take\textquotedblright{} (DEDR 2151): to take on agreement, to admit.
\end{cousinweb}

\begin{grammarlens}[title={What you've built}]
One more word for telling what happened.
\end{grammarlens}

\subsection*{Your turn}
\begin{itemize}
\raggedright
  \item \emph{Say it:} \emph{oppukkoḷ}
  \item \emph{Say it:} \emph{oppukkoḷ}, once more
  \item \emph{Say it:} \emph{kaṭaiciyil avaṉ taṉ kuṟṟattai oppukkoṇṭāṉ.} — In the end he confessed his crime.
  \item \emph{Recall:} say the Tamil for to cordon off, then the Tamil for a prisoner, then say \emph{oppukkoḷ} again
\end{itemize}

\begin{tcolorbox}[breakable,colback=teal!4,colframe=teal!35!black,title={Before you move on}]
What does \ta{ஒப்புக்கொள்} mean? (\textquotedblleft{}To confess; to admit\textquotedblright{}.) Say it once more.
\end{tcolorbox}
